% This LaTeX file was generated from the uploaded profile PDF.
\documentclass[11pt,a4paper]{article}

\usepackage[margin=1in]{geometry}
\usepackage[T1]{fontenc}
\usepackage[utf8]{inputenc}
\usepackage{lmodern}
\usepackage[hidelinks]{hyperref}
\usepackage{enumitem}
\usepackage{titlesec}
\usepackage{xcolor}

\setlength{\parindent}{0pt}
\setlength{\parskip}{0.4em}
\pagestyle{empty}

\titleformat{\section}{\large\bfseries}{}{0em}{}
\titlespacing*{\section}{0pt}{0.8em}{0.4em}

\begin{document}

\begin{center}
    {\LARGE \textbf{Mohammed Bilal p s}}\\[0.4em]
    PhD candidate @ ICB (Quantum Control) \;|\; IISER BPR (Physics) \;|\; IITM (Data Science) \\
    Chanakya Fellowship (I-Hub Quantum Technology) recipient (PG-level) -- 2024\\[0.6em]
    France\\
    \href{mailto:billabobz.official@gmail.com}{billabobz.official@gmail.com} \;|\;
    \href{https://www.linkedin.com/in/mohammedbilalps}{linkedin.com/in/mohammedbilalps}
\end{center}

\section*{Summary}
PhD candidate at the University of Burgundy (ICB Dijon), working under Dominique Sugny and Camille on quantum control theory with a focus on Pontryagin-based optimal control methods. Research interests include numerical and gradient-based optimization techniques for solving quantum control problems, with particular emphasis on robustness under realistic constraints and noise. Especially interested in integrating machine learning approaches to enhance control strategies and improve scalability for practical quantum technologies. Beyond control problems, actively exploring quantum error correction, fault-tolerant architectures, and data-driven methods relevant to the quantum industry. Values intellectual flexibility and interdisciplinary thinking, viewing research as a continuous balance between exploration and exploitation: testing new ideas while refining what works.

\section*{Top Skills}
\begin{itemize}[leftmargin=1.5em]
    \item Quantum control
    \item Quantum information
    \item Condensed matter physics
    \item Machine larning and ai methods in physics
    \item RL
    \item numerical mehtods
\end{itemize}

\section*{Experience}
\textbf{IISER Berhampur} \hfill \textit{May 2023 -- May 2024}\\
\textit{MS Thesis} \hfill India
\begin{itemize}[leftmargin=1.5em]
    \item Worked on ``optimal charging protocol for multi-qubit quantum battery'' under Dr. Victor Mukherjee and Paolo Andrea Erdman.
\end{itemize}

\textbf{revamp24} \hfill \textit{June 2020 -- May 2024}\\
\textit{Math Teacher} \hfill Bangalore

\textbf{Six Red Marbles} \hfill \textit{2019 -- 2021}\\
\textit{Mathematics Specialist} \hfill Chennai, Tamil Nadu, India
\begin{itemize}[leftmargin=1.5em]
    \item Subject Matter Expert in math and calculus.
\end{itemize}

\section*{Education}
\textbf{Universit\'e de Bourgogne} \hfill \textit{November 2024 -- November 2027}\\
Doctor of Philosophy (PhD), Theoretical and Mathematical Physics

\textbf{Indian Institute of Science Education and Research (IISER), Berhampur} \hfill \textit{2019 -- 2024}\\
BS--MS Integrated, Physics

\textbf{Indian Institute of Technology, Madras} \hfill \textit{September 2021 -- September 2026}\\
Bachelor of Science (BS), Data Science

\textbf{National Huda Central School}\\
Higher Secondary Education, Computer-Math
\section{works:}
\begin{itemize}
    \item quantm battery chrgining using rl for spin chain model-chanakya fellowship recpined masters work
    \item quantum contorl
\end{itemize}
\end{document}

